% -----------------------------------------------------------------------
% CCN Conference - Extended Abstract Track Template
% For papers up to 2 pages (excluding references); no supplement.
% -----------------------------------------------------------------------
\documentclass{ccn}  % options: submission (default), preprint, extended-abstract
\usepackage{macros}

\addbibresource{references.bib}

\title{Learning long-range spatial dependencies with spectral convolutional state space models}

\author{%
  Author One\affmark{1} \And
  Author Two\affmark{1} \And
  Author Three\affmark{2} \And
  Author Four\affmark{2,3} \AND
  Author Five\affmark{3} \And
  Author Six\affmark{4} \And
  Author Seven\affmark{4} \And
  Author Eight\affmark{1,4}%
}
\affiliation{1}{Department of Neuroscience, First University}
\affiliation{2}{Institute for Brain Research, Second University}
\affiliation{3}{Center for Neural Computation, Third University}
\affiliation{4}{Department of Psychology, Fourth University}
\emails{author1@first.edu, author3@second.edu, corresponding@fourth.edu}


\begin{document}

\maketitle

\begin{abstract}
Visual perception emerges from computations performed over feedforward and recurrent feedback connections between neurons. Yet today's state-of-the-art deep neural networks are almost entirely feedforward. Scaling has favored exclusively feedforward architectures that fit deep learning hardware, even at the cost of biological alignment. Here, we bridge the gap between biology and deep learning at scale with spectral convolutional state space models (SCSSMs), a scalable vision architecture that implements recurrent horizontal connections by lifting convolutional recurrence into the Fourier domain. A single SCSSM layer achieves human-level performance on visual reasoning challenges where transformers fail, and can be scaled up to large-scale vision benchmarks while matching transformer training speeds with fewer parameters and improved performance.
\end{abstract}

Recurrent neural networks are a central concept for explaining brains and behavior. In vision, local recurrent circuits through horizontal connections support contour integration, perceptual grouping, and figure-ground segmentation~\citep{gilbert1989,bosking1997,field1993,gilbertli2013,mely2018}. The horizontal gated recurrent unit (hGRU)~\citep{linsley2018} demonstrated that these connections could be incorporated into a deep learning architecture and trained end-to-end to solve multiple perceptual challenges~\citep{linsley2020}. But like all nonlinear RNNs, the hGRU is inherently sequential, rendering training slow and poorly suited for the capabilities of GPUs and TPUs. In contrast, Transformers~\citep{vaswani2017} are an exquisite fit for GPU and TPU capabilities and scaling, but struggle on tasks that demand structured long-range spatial integration~\citep{tay2021} and have are often misaligned with human perception~\citep{Linsley2025-hz}.

\paragraph{State space models} (SSMs)~\citep{gu2022} are a potential bridge between the architectural preferences of today's deep learning hardware and the architectural features of brains: linear recurrences that admit efficient parallel training. But extending SSMs to vision data has required strict constraints that make it impossible to implement biologically inspired recurrent connections. The S4ND~\citep{nguyen2022} separately computes dynamics along each axis of spatiotemporal data, while ConvS5~\citep{smith2023convs5} restricts the recurrence to 1$\times$1 convolutions. These choices reflect a structural barrier: composing spatial convolutions during a parallel scan causes kernel sizes to grow with each step, making the algorithm quadratic. We resolve this tension by lifting the convolutional SSM into the Fourier domain, where convolution composition becomes pointwise multiplication, restoring linear complexity while preserving full spatial coupling.

We call the resulting model the spectral convolutional SSM (SCSSM). We believe it is the first SSM that implements recurrent horizontal connections within a scalable and parallelizable architecture. On PathFinder and PathTracker, two tasks which require long-range spatial integration over static and dynamic stimuli, a one-layer SCSSM approaches human-level accuracy with fewer parameters and similar training speeds to Transformers (\Cref{Zoo-Results}). Moreover, we find early positive evidence for scaling SCSSMs: swapping the Transformers in a state-of-the-art edge vision architecture (the ShViT) with SCSSMs leads to better performance in a model with around 5\% fewer parameters (74.98\% vs. 73.66\% accuracy).

\begin{figure*}[!ht]\vspace{-2mm}
\includegraphics[width=\textwidth]{figures/fig1.pdf}
\caption{\textbf{A single layer of the SCSSM rivals humans on visual reasoning tasks.} We compare the SCSSM, leading SSMs, Transformers, and humans on two visual reasoning challenges: \textbf{(a)}~PathFinder, and \textbf{(b)}~PathTracker. \textbf{(c--d)}~SCSSMs and its gated variants are more accurate and efficient than leading sequence models (Transformers and SSMs) and other recurrent convolution vision models. Removing the spectral formulation (No-FFTControl) reduces performance to chance, confirming that the spectral formulation is essential.}
\label{Zoo-Results}
\end{figure*}\vspace{-2mm}

\paragraph{The Spectral Convolutional SSM}
An SSM describes the continuous-time evolution of a hidden state $\hvec(t)$ (\Cref{eq:ssm}), which can be discretized as a recurrence (\Cref{eq:ssm:recurrence}):
 
\begin{minipage}[t]{.50\linewidth}
\vspace{-4mm}
  \begin{subequations}
    \label{eq:ssm}
    \begin{align}
      \hvec'(t) &= \Amat \hvec(t) + \Bmat \xvec(t) \\
      \yvec(t) &= \Cmat \hvec(t) + \Dmat \xvec(t)
    \end{align}
  \end{subequations}
\end{minipage}%
\begin{minipage}[t]{.45\linewidth}
\vspace{-4mm}
  \begin{subequations}
    \label{eq:ssm:recurrence}
    \begin{align}
    \label{eq:ssm:recurrence:1}
      \hvec_{t} &= \overline{\Amat} \hvec_{t-1} + \overline{\Bmat} \xvec_t \\
    \label{eq:ssm:recurrence:2}
      \yvec_t &= \Cmat \hvec_t + \Dmat \xvec_t
    \end{align}
  \end{subequations}
\end{minipage}%
 
\paragraph{Convolutional SSM.} To model horizontal connectivity, we replace the matrix products in \Cref{eq:ssm} with spatial convolutions:
%
\begin{subequations}
    \label{eq:conv_ssm_continuous}
    \begin{align}
    \label{eq:conv_ssm_continuous:1}
      \frac{\partial \hvec}{\partial t}(r,t) &= \Kmat_h \ast \hvec(r, t) + \Kmat_x \ast \xvec(r, t) \\
      \yvec(r, t) &= \Kmat_y \ast \hvec(r, t)
    \end{align}
\end{subequations}
%
where $r$ indexes spatial position and $\ast$ denotes convolution. This is not scalable: composing $\Kmat_h$ during a parallel scan makes complexity quadratic~\citep{smith2023convs5}.
 
\paragraph{Spectral formulation.} In the Fourier domain, convolution becomes pointwise multiplication. We apply the 2D FFT to the spatial dimensions, yielding:
%
\begin{align}
\label{eq:convssm}
\Amat = \Lmat, \qquad
\Bmat = \text{FFT} \circ \Kmat_x, \qquad
\Cmat = \Kmat_y \circ \text{iFFT}
\end{align}
%
where $\Lmat$ is a learnable diagonal in frequency space. The parallel scan thus reduces to scalar operations at each frequency, restoring linear complexity with full spatio-temporal coupling. On top of this core SCSSM, we also build versions with Mamba-style (Mamba-SCSSM) and Gated Delta Network-style gating (GDN-SCSSM).

\paragraph{Experiments} We evaluate SCSSMs on \textbf{PathFinder} \citep{linsley2018}, a static contour-tracing task (\Cref{Zoo-Results}a), and \textbf{PathTracker}, a dynamic 32-frame dot-tracking task (\Cref{Zoo-Results}b). We compare against Transformers (with spatial or spatiotemporal attention), sequence-based SSMs (Mamba2, GDN), and convolutional SSMs without the spectral formulation (Mamba-CSSM, S5-CSSM).

\paragraph{Results.} SCSSM variants approach human accuracy on both tasks (\Cref{Zoo-Results}c--d). In contrast, non-spectral convolutional SSMs were at chance (${\sim}50\%$), confirming that the kernel growth problem prevents learning entirely. Transformers also struggle, while sequence-based SSMs are better, but still lag behind the SCSSM.

\printbibliography

\end{document}
